\icode{N13L-MG-SI-L2-261019-001}
\nceeID[5]{2013山东}{7}
双星系统由两颗恒星组成，两恒星在相互引力的作用下，分别围绕其连线上的某一点做周期相同的匀速圆周运动。研究发现，双星系统演化过程中，两星的总质量、距离和周期均可能发生变化。若某双星系统中两星做圆周运动的周期为\(T\)，经过一段时间演化后，两星总质量变为原来的\(k\)倍，两星之间的距离变为原来的\(n\)倍，则此时圆周运动的周期为\paren[]
\begin{choices}
	\item \(\sqrt{\dfrac{n}{k}}T\)
	\item \(\sqrt{\dfrac{n^2}{k}}T\)
	\item \(\sqrt{\dfrac{n^3}{k^2}}T\)
	\item \(\sqrt{\dfrac{n^3}{k}}T\)
\end{choices}
\begin{answer}
	D\qquad 提示：\(\frac{GMm}{L^2} = m \frac{4\pi^2}{T^2}r_1\)、\(\frac{GMm}{L^2} = M \frac{4\pi^2}{T^2}r_2\)、\(r_1 + r_2 = L\)联立：\(\frac{L^3}{T^2} = \frac{G(M + m)}{4\pi^2} \)
\end{answer}